% !TEX root = ../../hw07.tex
% Source: author's supplied text in chat, 2026-11-06, Q5.
% Preserve factorisation, both sign arguments, and maximum-value calculation.
\begin{proof}[Solution]
  Let
  \[
    f(x)=\frac1{1+x}-\frac1{1+ax},
    \qquad x\ge0,\quad a>1.
  \]
  To find the global maximum $x^*$ such that $f(x)\le f(x^*)$,
  first consider $f'(x^*)=0$,
  \[
    \begin{aligned}
      f'(x)&=\frac{a}{(1+ax)^2}-\frac1{(1+x)^2}\\[6pt]
        &=\frac{a(1+x)^2-(1+ax)^2}{(1+ax)^2(1+x)^2}\\[6pt]
        &=\frac{(a-1)(1-ax^2)}{(1+ax)^2(1+x)^2}.
    \end{aligned}
  \]
  Since $a>1$ and the denominator is strictly positive for $x\ge0$,
  the critical points satisfy
  \[
    \begin{aligned}
      f'(x)=0&\iff 1-ax^2=0\\
        &\iff x^2=\frac1a\\
        &\iff x=\pm\frac1{\sqrt a}\\
        &\iff x=\frac1{\sqrt a},\qquad x\ge0.
    \end{aligned}
  \]
  The sign of $f'(x)$ is determined by the $(1-ax^2)$ term.
  For $x<1/\sqrt a$,
  \[
    \begin{aligned}
      0\le x<\frac1{\sqrt a}
        &\implies 0\le x^2<\frac1a\\
        &\implies ax^2<1\\
        &\implies 1-ax^2>0\\
        &\implies f'(x)>0.
    \end{aligned}
  \]
  Similarly, for $x>1/\sqrt a$,
  \[
    \begin{aligned}
      0<\frac1{\sqrt a}<x
        &\implies 0<\frac1a<x^2\\
        &\implies 1<ax^2\\
        &\implies 1-ax^2<0\\
        &\implies f'(x)<0.
    \end{aligned}
  \]
  Since $f'(x)>0$ for $0\le x<1/\sqrt a$ and $f'(x)<0$ for
  $x>1/\sqrt a$, the function $f$ is strictly increasing before
  $1/\sqrt a$ and strictly decreasing afterward.
  This follows from the mean value theorem on
  any closed interval on either side of $1/\sqrt a$, including
  intervals starting at $0$. The sign change is illustrated in
  \autoref{fig:hw07:rational-maximum}.
  \begin{marginfigure}
    \centering
    \input{\analysisroot/figures/interval-styles.tex}
\begin{tikzpicture}[interval diagram,x=1.45cm,y=6.5cm]
  \draw[->] (0,0)--(2.35,0) node[right=4pt] {$x$};
  \draw[->] (0,0)--(0,.43) node[above=4pt] {$y$};
  \node[below left=5pt] at (0,0) {$0$};
  \draw[gray,dashed] (0,1/3)--(2.17,1/3);
  \node[left=5pt] at (0,1/3) {$1/3$};
  \draw[gray,dotted] (.5,0)--(.5,1/3);
  \draw (.5,0)--(.5,-.008) node[below=5pt] {$1/2$};
  \draw[->,domain=0:2.17,samples=90]
    plot (\x,{1/(1+\x)-1/(1+4*\x)});
  \fill (.5,1/3) circle[radius=1.6pt];
  \node[above right=6pt] at (.5,1/3) {unique maximum};
  \node at (.23,-.13) {$f'>0$};
  \node at (1.52,-.13) {$f'<0$};
  \draw[->] (.08,-.195)--(.4,-.165);
  \draw[->] (.75,-.165)--(2.08,-.195);
\end{tikzpicture}

    \caption[A unique maximum from a derivative sign change.]{%
      The case $a=4$. The function increases up
      to $x^*=1/2$ and decreases afterward. Its maximum is $1/3$.
      The sign argument in the proof applies to every $a>1$.}
    \label{fig:hw07:rational-maximum}
  \end{marginfigure}
  Therefore, $f$ attains its unique global maximum at
  $x^*=1/\sqrt a$,
  \[
    \begin{aligned}
      f(x)&\le f\left(\frac1{\sqrt a}\right)\\[6pt]
        &=\frac1{1+\frac1{\sqrt a}}-
          \frac1{1+\frac{a}{\sqrt a}}\\[6pt]
        &=\frac{\sqrt a}{\sqrt a+1}-
          \frac{\sqrt a}{\sqrt a+a}\\[6pt]
        &=\frac{\sqrt a}{\sqrt a+1}-\frac1{1+\sqrt a}\\[6pt]
        &=\frac{\sqrt a-1}{\sqrt a+1}.
    \end{aligned}
  \]
  Equality holds if and only if $x=1/\sqrt a$.
\end{proof}
